\documentclass[a4paper,12pt]{article}

% --- Pakiety ---
\usepackage[polish]{babel}
\usepackage[utf8]{inputenc}
\usepackage[T1]{fontenc}
\usepackage[top=2.5cm,bottom=2.5cm,left=3cm,right=3cm]{geometry}
\usepackage{amsmath}
\usepackage{amssymb}
\usepackage{graphicx}
\usepackage{booktabs}
\usepackage{array}
\usepackage{float}
\usepackage[colorlinks=true, allcolors=blue]{hyperref}
\usepackage{titlesec}

\titlespacing*{\subsection}{0pt}{12pt}{6pt}
\setlength{\parindent}{15pt}
\setlength{\parskip}{6pt}
\renewcommand{\arraystretch}{1.15}
\setlength{\tabcolsep}{3pt}

% --- Tytuł ---
\title{
    \vspace{1.5cm}
    \textbf{PODSTAWY TEORII INFORMACJI}\\
    \vspace{1.5cm}
    Laboratorium 5\\
    Eksperyment pragmatyczny i wnioski końcowe\\
    \vspace{1.5cm}
    \huge \textbf{Użyteczność prognoz meteorologicznych w decyzjach użytkowników}\\
    \vspace{1cm}
}

\author{
Maciek Chotkowski \\ Martsin Lazouski \\ Nadia Ziecina \\
Michał Kasznia \\ Kasiński Julian \\ Volodymyr Kamenchuk \\ Usevalad Nikalaichuk
}

\date{29 maja 2026}

\begin{document}

\maketitle
\newpage

\tableofcontents
\newpage

\section{Cel raportu}

Celem raportu jest sprawdzenie, czy nasze narzędzie pogodowe jest użyteczne dla
różnych użytkowników. W poprzednich raportach ocenialiśmy głównie zgodność
prognoz z danymi referencyjnymi. W tym etapie interesuje nas coś innego: czy
prognoza, nawet jeśli jest technicznie poprawna, rzeczywiście pomaga ludziom w
podjęciu decyzji.

Eksperyment ma dwie części. Najpierw, dnia \textbf{18 maja 2026}, pytamy grupę
20 osób, co chcą zrobić następnego dnia. Następnie zapisujemy \textbf{naszą
prognozę} dla dnia \textbf{19 maja 2026} i podejmujemy \textbf{naszą decyzję}:
czy dana aktywność jest możliwa, niemożliwa albo możliwa warunkowo. Druga część
polega na tym samym, ale dla prognozy z tygodniowym wyprzedzeniem: dnia
\textbf{18 maja 2026} pytamy o plany na \textbf{25 maja 2026}. W tej części
trafność powinna być mniejsza, ponieważ im dalszy horyzont prognozy, tym większa
niepewność pogody i samego użytkownika.

\section{Pytania ogólne}

Najpierw nie pytaliśmy o jedną konkretną aktywność, tylko o potrzeby
użytkowników. Użyliśmy prostego zestawu pytań:

\begin{enumerate}
    \item Co chcesz zrobić jutro na zewnątrz?
    \item O której godzinie planujesz tę aktywność?
    \item Które warunki są dla Ciebie najważniejsze: temperatura, opady, wiatr,
    stan podłoża, bezpieczeństwo, komfort?
    \item Przy jakich warunkach zrezygnujesz?
    \item Czy decyzja zależy tylko od Ciebie, czy także od kogoś innego, np.
    trenera, organizatora, rodzica albo administratora obiektu?
    \item Czy po wykonaniu lub odwołaniu planu oceniasz prognozę jako użyteczną?
\end{enumerate}

Takie pytania pozwalają przejść od samej prognozy do modelu użytkownika. Dla
jednej osoby lekki deszcz jest małym problemem, dla innej oznacza całkowitą
rezygnację. Dlatego ta sama prognoza może mieć różną wartość praktyczną.

\section{Daty, nasza prognoza i nasza decyzja}

Badanie zapisujemy w konkretnym harmonogramie, żeby później można było wrócić
do danych i sprawdzić, czy model rzeczywiście pomógł w decyzji.

\begin{table}[H]
\centering
\small
\begin{tabular}{p{3.0cm}p{10.0cm}}
\toprule
\textbf{Data} & \textbf{Czynność} \\
\midrule
18.05.2026 & pytamy 20 osób o plan na następny dzień i ich warunki rezygnacji \\
18.05.2026 wieczorem & zapisujemy naszą prognozę na 19.05.2026 i naszą decyzję dla każdej osoby \\
19.05.2026 & sprawdzamy, czy użytkownicy faktycznie mogli wykonać plan \\
25.05.2026 & sprawdzamy osobny eksperyment: prognozę wykonaną tydzień wcześniej \\
26.05.2026 & po tygodniu sprawdzamy wybrane pomiary: temperaturę, opady i wiatr \\
\bottomrule
\end{tabular}
\caption{Harmonogram eksperymentu.}
\end{table}

Nasza prognoza robocza dla Warszawy na 19 maja 2026 była następująca:

\begin{table}[H]
\centering
\small
\begin{tabular}{p{2.0cm}p{2.6cm}p{2.6cm}p{2.4cm}p{4.0cm}}
\toprule
\textbf{Godzina} & \textbf{Temperatura} & \textbf{Opady} & \textbf{Wiatr} & \textbf{Nasza decyzja ogólna} \\
\midrule
08:00 & około $12^\circ C$ & słaby opad & słaby & aktywności wrażliwe na mokrą nawierzchnię: warunkowo albo nie \\
12:00 & około $17^\circ C$ & brak opadów & słaby & większość aktywności możliwa \\
17:00 & około $19^\circ C$ & brak opadów & słaby & sport i spacer możliwe, ale stan podłoża trzeba sprawdzić osobno \\
20:00 & około $17^\circ C$ & brak opadów & słaby & aktywności krótkie możliwe, decyzja zależy od użytkownika \\
\bottomrule
\end{tabular}
\caption{Nasza prognoza i decyzja ogólna dla dnia 19.05.2026.}
\end{table}

Już na tym etapie widać problem pragmatyczny. Pogoda po południu wygląda
dobrze, więc model może powiedzieć: ,,gra albo spacer są możliwe''. Jednak stan
boiska, wcześniejsze opady, zdrowie zawodnika albo decyzja trenera nie wynikają
bezpośrednio z samej prognozy.

\section{Grupa badana}

W eksperymencie uwzględniono 20 osób z różnymi potrzebami. Nie chodziło o to,
aby uzyskać statystycznie pełną próbę, ale żeby pokazać różne typy sytuacji, w
których prognoza pogody jest wykorzystywana.

\begin{table}[H]
\centering
\scriptsize
\begin{tabular}{p{0.8cm}p{2.6cm}p{3.4cm}p{4.2cm}p{2.5cm}}
\toprule
\textbf{Nr} & \textbf{Grupa} & \textbf{Plan} & \textbf{Najważniejsze warunki} & \textbf{Kto decyduje?} \\
\midrule
1 & piłkarz & trening na boisku & opady, stan murawy, kontuzje & trener i zawodnik \\
2 & piłkarz & mecz amatorski & opady, wiatr, dostępność boiska & organizator \\
3 & biegacz & trening 10 km & temperatura, opady, wiatr & użytkownik \\
4 & rowerzysta & dojazd do pracy & opady, wiatr, bezpieczeństwo & użytkownik \\
5 & rodzic & spacer z dzieckiem & temperatura, opady, wiatr & rodzic \\
6 & uczeń & droga do szkoły & opady, temperatura & rodzic i uczeń \\
7 & student & przejście między zajęciami & opady, komfort & użytkownik \\
8 & pracownik budowy & praca na zewnątrz & deszcz, wiatr, bezpieczeństwo & kierownik \\
9 & ogrodnik & prace w ogrodzie & opady, wilgotność podłoża & użytkownik \\
10 & fotograf & zdjęcia plenerowe & opady, zachmurzenie, światło & klient i fotograf \\
11 & kierowca & dłuższa trasa & opady, śliskość, widoczność & użytkownik \\
12 & kurier & dostawy rowerowe & deszcz, wiatr, czas & firma i użytkownik \\
13 & trener & zajęcia sportowe & opady, bezpieczeństwo grupy & trener \\
14 & właściciel psa & spacer & opady, temperatura & użytkownik \\
15 & senior & wyjście do sklepu & śliskość, wiatr, temperatura & użytkownik \\
16 & organizator & małe wydarzenie plenerowe & opady, wiatr, komfort grupy & organizator \\
17 & turysta & krótka wycieczka & opady, wiatr, temperatura & grupa \\
18 & operator drona & lot testowy & wiatr, opady, widoczność & użytkownik \\
19 & pracownik kawiarni & ogródek zewnętrzny & opady, temperatura, wiatr & właściciel \\
20 & osoba po urazie & lekki trening & śliskość, opady, ryzyko kontuzji & użytkownik \\
\bottomrule
\end{tabular}
\caption{Przykładowa grupa 20 osób o różnych potrzebach pogodowych.}
\end{table}

\section{Przewidywanie i odpowiedzi po fakcie}

Dla każdej osoby zapisaliśmy datę aktywności, naszą decyzję oraz późniejszą
odpowiedź użytkownika. Odpowiedź ,,warunkowo'' oznaczała, że aktywność jest
możliwa, ale istnieje ryzyko albo trzeba zmienić plan.

\begin{table}[H]
\centering
\scriptsize
\begin{tabular}{p{0.8cm}p{2.0cm}p{2.7cm}p{2.0cm}p{2.0cm}p{4.0cm}}
\toprule
\textbf{Nr} & \textbf{Godz.} & \textbf{Plan} & \textbf{Nasza decyzja} & \textbf{Faktycznie} & \textbf{Interpretacja} \\
\midrule
1 & 17:00 & trening piłkarski & warunkowo & tak & trener zdecydował, że zajęcia się odbędą mimo mokrej murawy \\
2 & 18:00 & mecz amatorski & warunkowo & nie & administrator zamknął boisko po wcześniejszych opadach \\
3 & 12:00 & bieganie & tak & tak & lekki wiatr nie miał większego znaczenia \\
4 & 08:00 & rower do pracy & warunkowo & nie & użytkownik uznał mokrą nawierzchnię za zbyt ryzykowną \\
5 & 08:00 & spacer z dzieckiem & nie & nie & poranny opad był wystarczającym powodem rezygnacji \\
6 & 08:00 & droga do szkoły & tak & tak & prognoza wystarczyła do przygotowania ubioru \\
7 & 12:00 & przejście między zajęciami & tak & tak & aktywność była krótka, więc pogoda nie przeszkodziła \\
8 & 08:00 & praca budowlana & nie & nie & kierownik odwołał część pracy ze względów bezpieczeństwa \\
9 & 17:00 & ogród & warunkowo & nie & podłoże było zbyt mokre, czego model nie przewidział dokładnie \\
10 & 17:00 & zdjęcia & tak & nie & brak opadów nie wystarczył, bo problemem było światło i zachmurzenie \\
11 & 12:00 & trasa autem & tak & tak & warunki były akceptowalne \\
12 & 12:00 & dostawy rowerowe & warunkowo & tak & firma wymagała realizacji, użytkownik dostosował ubiór \\
13 & 17:00 & zajęcia sportowe & warunkowo & tak & trener zaakceptował ryzyko i zmienił intensywność zajęć \\
14 & 20:00 & spacer z psem & tak & tak & użytkownik zaakceptował lekko mokre warunki \\
15 & 08:00 & wyjście seniora & warunkowo & nie & osoba zrezygnowała przez obawę przed śliską nawierzchnią \\
16 & 17:00 & wydarzenie plenerowe & warunkowo & nie & organizator odwołał wydarzenie z powodu ryzyka dla grupy \\
17 & 12:00 & wycieczka & warunkowo & tak & grupa skróciła trasę \\
18 & 17:00 & lot dronem & tak & tak & wiatr był słaby, więc próba była możliwa \\
19 & 20:00 & ogródek kawiarni & warunkowo & nie & właściciel uznał, że będzie za mało klientów \\
20 & 17:00 & lekki trening po urazie & warunkowo & nie & indywidualne ryzyko kontuzji było ważniejsze niż sama pogoda \\
\bottomrule
\end{tabular}
\caption{Nasza decyzja i odpowiedź użytkownika dla dnia 19.05.2026.}
\end{table}

\section{Kontrola pomiarów po tygodniu}

Po tygodniu warto sprawdzić kilka konkretnych pomiarów, a nie tylko polegać na
odpowiedziach użytkowników. Dlatego dla dnia 19.05.2026 zapisujemy wybrane
godziny i parametry do późniejszej kontroli.

\begin{table}[H]
\centering
\small
\begin{tabular}{p{2.4cm}p{2.4cm}p{2.4cm}p{2.4cm}p{4.0cm}}
\toprule
\textbf{Godzina} & \textbf{Temperatura} & \textbf{Opady} & \textbf{Wiatr} & \textbf{Znaczenie dla decyzji} \\
\midrule
08:00 & $12.0^\circ C$ & 0.3 mm & 1.5 m/s & mokra nawierzchnia mogła zmienić decyzję rowerzysty i seniora \\
12:00 & $16.9^\circ C$ & 0.0 mm & 1.1 m/s & dobre warunki dla krótkich aktywności \\
17:00 & $18.6^\circ C$ & 0.0 mm & 1.7 m/s & pogoda dobra, ale stan boiska i organizacja nadal istotne \\
20:00 & $16.9^\circ C$ & 0.0 mm & 1.9 m/s & sama pogoda nie tłumaczy rezygnacji kawiarni \\
\bottomrule
\end{tabular}
\caption{Wybrane pomiary do kontroli po tygodniu.}
\end{table}

Ta kontrola jest ważna, bo oddziela dwa rodzaje błędów. Pierwszy to błąd
prognozy pogody, np. źle przewidziany deszcz. Drugi to błąd użyteczności, gdy
pogoda została przewidziana poprawnie, ale decyzja użytkownika zależała od
czegoś innego.

\section{Osobny eksperyment: prognoza za tydzień}

Drugi eksperyment wykonujemy dla dłuższego horyzontu czasowego. Dnia
\textbf{18.05.2026} pytamy te same osoby o plany na \textbf{25.05.2026}, czyli
na tydzień później. Nasza prognoza nie jest już traktowana tak pewnie jak
prognoza jednodniowa. Przy prognozie tygodniowej model może jeszcze wskazać
ogólny trend, ale trudniej przewidzieć konkretną godzinę opadu, stan boiska,
siłę wiatru i decyzje organizacyjne.

Nasza decyzja dla prognozy tygodniowej była więc ostrożniejsza. Częściej
używaliśmy odpowiedzi \textbf{warunkowo}, ponieważ przy takim horyzoncie nie da
się uczciwie obiecać, że warunki będą dokładnie takie jak w prognozie.

\begin{table}[H]
\centering
\scriptsize
\begin{tabular}{p{0.8cm}p{2.8cm}p{2.2cm}p{2.2cm}p{5.0cm}}
\toprule
\textbf{Nr} & \textbf{Plan na 25.05} & \textbf{Nasza decyzja 18.05} & \textbf{Faktycznie} & \textbf{Dlaczego wynik się zmienił?} \\
\midrule
1 & trening piłkarski & warunkowo & nie & boisko było zajęte, a trener zmienił plan zajęć \\
2 & mecz amatorski & warunkowo & tak & pogoda była gorsza niż oczekiwano, ale organizator nie odwołał meczu \\
3 & bieganie & tak & nie & użytkownik zrezygnował z powodów prywatnych, niezależnych od pogody \\
4 & rower do pracy & warunkowo & tak & opad przesunął się na inną godzinę niż przewidywano \\
5 & spacer z dzieckiem & warunkowo & nie & pojawił się silniejszy wiatr i rodzic uznał spacer za niekomfortowy \\
8 & praca budowlana & warunkowo & nie & decyzję podjął kierownik po sprawdzeniu warunków rano \\
10 & zdjęcia plenerowe & warunkowo & nie & problemem było zachmurzenie i światło, a nie sama temperatura \\
13 & zajęcia sportowe & warunkowo & tak & trener przeniósł zajęcia na późniejszą godzinę \\
16 & wydarzenie plenerowe & warunkowo & nie & organizator odwołał wydarzenie ze względu na ryzyko dla grupy \\
18 & lot dronem & nie & nie & w tym przypadku ostrożna decyzja była trafna, bo wiatr nadal był ryzykowny \\
\bottomrule
\end{tabular}
\caption{Przykładowe wyniki eksperymentu dla prognozy z tygodniowym wyprzedzeniem.}
\end{table}

W tym eksperymencie trafność jest wyraźnie mniejsza. Dla prognozy na następny
dzień większość decyzji była przynajmniej częściowo użyteczna, bo warunki
pogodowe nie zdążyły się mocno zmienić. Dla prognozy tygodniowej pojawia się
więcej błędów: zmienia się godzina opadów, użytkownik zmienia plan, trener
przesuwa zajęcia, organizator odwołuje wydarzenie albo stan boiska okazuje się
inny niż zakładaliśmy.

\begin{table}[H]
\centering
\small
\begin{tabular}{p{4.0cm}p{3.0cm}p{6.0cm}}
\toprule
\textbf{Horyzont prognozy} & \textbf{Trafność decyzji} & \textbf{Interpretacja} \\
\midrule
1 dzień & około 65--70\% & prognoza jest dość użyteczna, bo warunki i plany użytkowników są jeszcze stabilne \\
7 dni & około 40--50\% & trafność spada, bo rośnie niepewność pogody, organizacji i zachowania użytkowników \\
\bottomrule
\end{tabular}
\caption{Porównanie trafności dla prognozy jednodniowej i tygodniowej.}
\end{table}

Ten eksperyment pokazuje ważną rzecz: niższa trafność po tygodniu nie oznacza,
że narzędzie jest bezużyteczne. Oznacza tylko, że decyzja powinna być mniej
kategoryczna. Przy horyzoncie tygodniowym narzędzie powinno mówić raczej:
,,wstępnie wygląda dobrze'', ,,ryzyko jest podwyższone'' albo ,,sprawdź ponownie
dzień wcześniej''. Wtedy prognoza nadal wspiera planowanie, ale nie udaje
pewności, której w rzeczywistości nie ma.

\section{Interpretacja eksperymentu}

Wyniki obu eksperymentów pokazują, że prognoza może być teoretycznie poprawna,
ale jej użyteczność zależy od użytkownika oraz od horyzontu czasowego.
Najciekawsze są przypadki, w których kilka osób otrzymało podobną prognozę, ale
podjęło różne decyzje. Dla piłkarza trening mógł się odbyć, jeżeli trener
zaakceptował warunki. Dla innego zawodnika podobna pogoda nie wystarczyła,
ponieważ boisko zostało zamknięte.

To oznacza, że model pogodowy nie odpowiada wprost na pytanie: ,,czy na pewno
zrobię to jutro?''. Odpowiada raczej: ,,czy pogoda wygląda na wystarczającą do
tej aktywności?''. Między tymi pytaniami istnieje duża różnica.

Decyzję można zapisać jako:
\[
D = f(W, U, O, L)
\]
gdzie $W$ oznacza pogodę, $U$ użytkownika, $O$ organizację, a $L$ czynniki
losowe. Nasze narzędzie najlepiej przewiduje $W$, częściowo może uwzględnić
$U$, ale nigdy nie przewidzi w pełni $O$ i $L$.

\section{Wnioski z całego projektu}

Raport 1 pokazał, że różne modele pogodowe potrafią dawać różne wyniki i trzeba
je porównywać za pomocą metryk błędu. Raport 2 wprowadził model użytkownika i
pokazał, że informacja jest wartościowa dopiero wtedy, gdy wspiera decyzję.
Raport 3 zbudował techniczną część projektu: pobieranie danych, wspólny format
CSV i porównywanie źródeł. Raport 4 przesunął projekt w stronę własnego modelu
i interpretacji wyników.

Raport 5 pokazuje ograniczenie całego podejścia. Nawet jeśli prognoza jest
poprawna, użyteczność narzędzia zależy od wielu rzeczy poza pogodą: zdrowia,
obowiązków, decyzji trenera, decyzji administratora, stanu boiska, decyzji
organizatora albo indywidualnej tolerancji na ryzyko. Osobny eksperyment z
prognozą tygodniową pokazuje dodatkowo, że wraz z wydłużeniem horyzontu spada
trafność decyzji.

Końcowy wniosek jest taki, że nasze narzędzie może wspierać decyzje, ale nie
może ich całkowicie zastąpić. Prognoza jest jednym z elementów decyzji, ale nie
całą decyzją. Dlatego najlepszy system nie powinien tylko mówić ,,tak'' albo
,,nie'', ale także pokazywać ryzyko, warunki i możliwe ograniczenia. Właśnie w
tym sensie użyteczność prognozy zależy od użytkownika.

\end{document}
